\documentclass[10pt]{beamer}
\usetheme{metropolis} % Requer o tema metropolis instalado, ou troque por 'Madrid'
\usepackage[utf8]{inputenc}
\usepackage[portuguese]{babel}
\usepackage{booktabs}
\usepackage{amsmath}

\title{Implementação de Pós-Graduação}
\subtitle{Status do Projeto, Viabilidade Financeira e Tomada de Decisão}
\author{}
\date{\today}

\begin{document}

\maketitle

% Slide 1: Status de Implementação
\begin{frame}{Status de Implementação}
    \begin{block}{Bloco 1: Infraestrutura (Concluído)}
        \begin{itemize}
            \item \textbf{Hospedagem Live:} Vimeo ✅
            \item \textbf{Pagamentos:} Hotmart ✅
            \item \textbf{Interface do Aluno:} Área de Membros ✅
            \item \textbf{Emissão de Notas:} Hotmart ✅
        \end{itemize}
    \end{block}

    \begin{block}{Bloco 2: Integração (Concluído)}
        \begin{itemize}
            \item Conexão simultânea com a plataforma \textbf{Focus} ✅
        \end{itemize}
    \end{block}
\end{frame}

% Slide 2: O Desafio do Parcelamento
\begin{frame}{O Desafio Estratégico}
    \begin{alertblock}{Requisito Pendente}
        \textbf{Parcelamento em 24x:} Não disponível nas integrações atuais. ❌
    \end{alertblock}

    \vspace{0.5cm}
    \textbf{Impacto:} 
    A ausência do parcelamento estendido exige uma reavaliação imediata da estratégia de precificação ou do modelo tecnológico para evitar barreiras de conversão.
\end{frame}

% Slide 3: Caminhos Estratégicos
\begin{frame}{Caminhos para Decisão}
    \begin{table}[]
        \centering
        \small
        \begin{tabular}{lp{3cm}p{3.5cm}}
            \toprule
            \textbf{Opção} & \textbf{Estratégia} & \textbf{Avaliação} \\ \midrule
            \textbf{A} & 12x R\$ 600,00 & Mais rápida; resolve o problema de conversão. \\ \midrule
            \textbf{B} & Sistema Próprio & Flexível; exige alto investimento e fluxo de caixa. \\ \midrule
            \textbf{C} & 12x R\$ 1.200,00 & Alto ticket; elevado risco de não vender. \\ \bottomrule
        \end{tabular}
    \end{table}
\end{frame}

% Slide 4: Análise Financeira (Break-even)
\begin{frame}{Análise Financeira: Ponto de Equilíbrio}
    \textbf{Custos Fixos (Docentes):}
    \begin{equation*}
        300h \times R\$ 200,00/h = R\$ 60.000,00
    \end{equation*}

    \textbf{Custos Variáveis:}
    \begin{itemize}
        \item Taxa administrativa por aluno: $R\$ 500,00$
    \end{itemize}

    \textbf{Meta para Iniciar Turma (Ticket R\$ 7.200):}
    \begin{itemize}
        \item Para 9 alunos: $R\$ 64.500,00$ de custo total.
        \item \textbf{Ponto de Equilíbrio:} 10 alunos ($R\$ 72.000,00$).
    \end{itemize}
\end{frame}

% Slide 5: Projeção de Escala
\begin{frame}{Projeção de Escala (30 Alunos)}
    Considerando uma turma com 30 alunos matriculados:

    \begin{itemize}
        \item \textbf{Faturamento Bruto:} $R\$ 216.000,00$
        \item \textbf{Impostos (10\%):} $- R\$ 21.600,00$
        \item \textbf{Faturamento Líquido Apurado:} \textbf{R\$ 194.400,00}
    \end{itemize}

    \vfill
    \begin{center}
        \textbf{Resultado:} Margem operacional saudável para o projeto.
    \end{center}
\end{frame}

\end{document}